\exercisegroup

\begin{enumerate}
\def\labelenumi{\arabic{enumi})}
\item
  Let E be a vector space and A a convex symmetric convex subset of E. Let T, T' be two locally convex topologies on E and U, U' be the uniform structures defined by T, T' on E. In order that the uniform structure induced on A by U' should be finer than that induced by U, it is necessary and sufficient that every neighbourhood of 0 for the topology induced on A by T should be a neighbourhood of 0 for the topology induced on A by T'.
\item
  \begin{enumerate}
  \def\labelenumii{\alph{enumii})}
  \tightlist
  \item
    Give an example of a non-compact closed set in \(\mathbf{R}^2\) , whose convex envelope is not closed. b) Show that, in \(\mathbf{R}^n\) , the convex envelope of a compact set is compact (cf.~II, p.~66, exerc. 9, a)).
  \end{enumerate}
\end{enumerate}

¶ 3) Let I be the compact interval {[}0, 1{]} of R and F be the vector space \(\mathcal{C}(\mathrm{I},\mathbf{R})\) of continuous real valued functions defined in I. Let E be the product space \(\mathbf{R}^{\mathbf{F}}\) ; for all \(a\in\mathrm{I}\) , let \(\varepsilon_{a}\) be the element of E such that \(\varepsilon_{a}(f)=f(a)\) for all \(f\in\mathrm{F}\) .

\begin{enumerate}
\def\labelenumi{\alph{enumi})}
\item
  Show that, when \(x\) varies in I, the set \(\mathbf{K}\) formed by the \(\varepsilon_x\) is compact in E.
\item
  Let \(\lambda\) be an element of \(E\) such that \(\lambda(f) = \int_0^1 f(t)dt\) for all \(f\in F\) (Lebesgue measure).
\end{enumerate}

Show that, in E, \(\lambda\) belongs to the closure of the convex envelope of K but does not belong to this convex envelope (cf.~FVR, II, p.~7, prop. 5).

\begin{enumerate}
\def\labelenumi{\arabic{enumi})}
\setcounter{enumi}{3}
\tightlist
\item
  With the notations of II, p.~72, exerc. 1 suppose also that the space \(\mathbf{E}\) is locally convex.
\end{enumerate}

\begin{enumerate}
\def\labelenumi{\alph{enumi})}
\item
  There exists a positive continuous linear form \(g\) in \(\mathbf{E}\) that extends \(f\), if and only if \(f\) is bounded below in \(M \cap (\mathbf{W} + \mathbf{P})\) for at least one neighbourhood \(\mathbf{W}\) of 0 in \(\mathbf{E}\).
\end{enumerate}

\begin{enumerate}
\def\labelenumi{\alph{enumi})}
\setcounter{enumi}{1}
\tightlist
\item
  Given a point \(x \in \mathbf{E}\) , there exists a positive continuous linear form \(g\) in \(\mathbf{E}\) such that \(g(x) = 1\) , if, and only if, \(-x \notin \overline{\mathbf{P}}\) .
\end{enumerate}

\begin{enumerate}
\def\labelenumi{\arabic{enumi})}
\setcounter{enumi}{4}
\tightlist
\item
  \begin{enumerate}
  \def\labelenumii{\alph{enumii})}
  \tightlist
  \item
    Let E be an infinite dimensional normed space and T be its topology. Show that there exists on E a normed space topology \(T'\) that is strictly finer than that of T and a normed space topology \(T''\) that is strictly coarser than that of T (define the neighbourhoods of 0 for these topologies, using a basis of E put in the form \((a_{\alpha,n})\) where \(\alpha\) varies in an infinite set of indices A and n in the set of integers \(\geqslant 0\) and where \(\|a_{\alpha,n}\| = 1\) for the given norm on E).
  \end{enumerate}
\end{enumerate}

\begin{enumerate}
\def\labelenumi{\alph{enumi})}
\setcounter{enumi}{1}
\item
  Let \(p\) be the norm defining the topology \(\mathcal{T}'\) . Show that, if \(E\) is complete for the topology \(\mathcal{T}\) , then \(p\) cannot be lower semi-continuous in \(E\) for the topology \(\mathcal{T}\) (use Baire's th. cf.~III, p.~25, corollary). Deduce that the convex set \(A\) defined by the relation \(p(x) < 1\) does not contain any interior point for \(\mathcal{T}\) even though all its points are internal.
\item
  Deduce from \(b\) ), that, if \(E\) is complete for the topology \(\mathcal{T}\) , then there exists in \(E\) convex sets which are not closed for \(\mathcal{T}\) , of which the intersection with every linear variety of finite dimension is closed for \(\mathcal{T}\) (cf.~II, p.~67, exerc. 12).
\end{enumerate}

\begin{enumerate}
\def\labelenumi{\arabic{enumi})}
\setcounter{enumi}{5}
\tightlist
\item
  Let \(\mathbf{E}\) be a vector space with its finest locally convex topology.
\end{enumerate}

\begin{enumerate}
\def\labelenumi{\alph{enumi})}
\item
  Show that every vector subspace of E is closed, and that, if M, N are two subspaces that are vectorial complements in E, then E is the direct topological sum of M and N. If \((e_{t})_{t\in I}\) is a basis for E, then E is the direct topological sum of the subspaces \(Re_{t}\) .
\item
  Let \(\mathrm{F}\) be a locally convex space whose topology is also the finest locally convex topology. Show that every linear mapping of \(\mathrm{E}\) in \(\mathrm{F}\) is a strict morphism.
\end{enumerate}

\begin{enumerate}
\def\labelenumi{\arabic{enumi})}
\setcounter{enumi}{6}
\item
  \begin{enumerate}
  \def\labelenumii{\alph{enumii})}
  \tightlist
  \item
    Let A be a convex set with at least one interior point in a topological vector space E. Show that the set of internal points of A is identical with the interior of A (cf.~exerc. 5, b)). b) Show that in the normed space \(E = l^{1}(\mathbf{N})\) , the convex cone P defined in II, p.~66, exerc. 11, b), generates E but does not contain any internal point.
  \end{enumerate}
\item
  Let E be a vector space with an enumerable basis and with the finest locally convex topology. Show that, if A is a set in E whose intersection with every vector subspace of finite dimension is closed in E, then A is closed in E (cf.~exerc. 5, c)).
\item
  Let E and F be two vector spaces each with its finest locally convex topology.
\end{enumerate}

\begin{enumerate}
\def\labelenumi{\alph{enumi})}
\item
  Show that if E and F each have an enumerable basis then every bilinear mapping of \(E \times F\) in a locally convex space G is continuous (use Du Bois-Reymond's th. FVR, V, p.~53, exerc. 8).
\item
  If one of the spaces E, F has a basis with cardinal equal to that of the continuum, show that there exists a non-continuous bilinear form in \(E \times F\) . (Reduce to the case where \(E = R^{(N)}\) , \(F = R^{N}\) , so that F can be identified with \(E^{*}\) and the bilinear forms on \(E \times F\) correspond bijectively with the linear mappings of \(E^{*}\) in itself; then consider the identity mapping of \(E^{*}\) , and note that in \(R^{N}\) , a compact set for the product topology cannot be absorbent.)
\end{enumerate}

\begin{enumerate}
\def\labelenumi{\arabic{enumi})}
\setcounter{enumi}{9}
\item
  Let \((\mathrm{E}_n)\) be an infinite sequence of locally convex spaces and let E be the topological direct sum of the family \((\mathrm{E}_n)\) . Show that the topology of E is identical with the topology \(\mathcal{T}_0\) defined in I, p.~24, exerc. 14.
\item
  Let \(\mathbf{I}\) be an infinite non-enumerable set. On the vector space \(\mathrm{E} = \mathbf{R}^{(I)}\) , show that the finest locally convex topology is distinct from the topology \(\mathcal{T}_0\) defined in \(\mathbf{I}\) , p.~24, exerc. 14; for this prove that the set of the \(x = (\xi_{i})_{i\in I} \in \mathrm{E}\) such that \(|\sum_{i\in I} \xi_i| < 1\) , is open in \(\mathcal{T}\) but not in \(\mathcal{T}_0\) .
\item
  Let \(\mathbf{E}\) be a vector space with an enumerable basis \((e_n)\) . Let \(\mathbf{V}\) be the balanced convex envelope of the set of the \(e_n\) and let \(\mathbf{W}\) be the balanced convex envelope of the set of points
\end{enumerate}

\[
a _ {n} = e _ {n} + (n - 1) e _ {1} \quad (n \geqslant 1).
\]

Let \(T_{1}\) (resp. \(T_{2}\) ) be the locally convex topology on E for which a fundamental system of neighbourhoods of 0 is formed by the \(\lambda V\) (resp. \(\lambda W\) ) for \(\lambda > 0\) . Show that \(T_{1}\) and \(T_{2}\) are Hausdorff, but that the lower bound of \(T_{1}\) and \(T_{2}\) in the set of locally convex topologies on E is not Hausdorff (cf.~II, p.~80, exerc. 26).

\begin{enumerate}
\def\labelenumi{\arabic{enumi})}
\setcounter{enumi}{12}
\item
  With the hypotheses of II, §6, show that E is complete for topology T which is the inductive limit of the \(T_{n}\) , if and only if, for each integer n and every Cauchy filter F on \(E_{n}\) for the topology induced by T, there exists \(p \geqslant n\) such that F is convergent in \(E_{p}\) for the topology \(T_{p}\) .
\item
  Let E be the strict inductive limit of an increasing sequence of locally convex spaces \(E_{n}\) (II, p.~33). Show that the topology of E is the finest of the topologies compatible with the vector space structure of E, whether locally convex or not, and inducing on \(E_{n}\) a coarser topology than the given topology \(T_{n}\) . (Let \(V_{0}\) be a neighbourhood of 0 for such a topology T and \((\mathrm{V}_{n})_{n\geqslant0}\) , a sequence of neighbourhoods of 0 for T such that \(V_{n+1} + V_{n+1} \subset V_{n}\) for all \(n \geqslant 0\) ; for all \(n \geqslant 1\) , consider, in \(E_{n}\) , a convex neighbourhood \(W_{n}\) of 0 that is contained in \(E_{n} \cap V_{n}\) , and take the convex envelope of the union of the \(W_{n}\) in E.)
\item
  Let I be an infinite non-enumerable set. Let \(\mathfrak{F}(\mathrm{I})\) be the family of finite subsets of I and E the direct sum space \(\mathbf{R}^{(I)}\) . For every \(J \in \mathfrak{F}(\mathrm{I})\) , let \(F_{J}\) be the subspace \(\mathbf{R}^{\mathrm{J}}\) of E, the product of the factors whose indices belong to J, with the product topology; let \(g_{J}\) be the canonical injection of \(F_{J}\) in E. Show that there exists a topology \(\mathcal{T}_0\) on \(E\) that is compatible with the vector space structure of \(E\) , which make the \(g_{J}\) continuous and which is strictly finer than the finest locally convex topology \(\mathcal{T}\) which makes the \(g_{J}\) continuous. (Note that the set \(V\) of the \(x = (\xi_{t})_{t\in I}\) in \(E\) such that \(\sum_{i\in I}|\xi_i|^{1 / 2}\leqslant 1\) is a neighbourhood of 0 for a topology compatible with the vector
\end{enumerate}

space structure of E, not containing any absorbent, symmetric, convex set.)

\begin{enumerate}
\def\labelenumi{\arabic{enumi})}
\setcounter{enumi}{15}
\tightlist
\item
  \begin{enumerate}
  \def\labelenumii{\alph{enumii})}
  \tightlist
  \item
    Let E be a vector space with an enumerable basis. Show that the finest locally convex topology on E is the finest of the topologies on E (compatible or not with the vector space structure of E) which induces the canonical topology on every finite dimensional subspace of E.
  \end{enumerate}
\end{enumerate}

\begin{enumerate}
\def\labelenumi{\alph{enumi})}
\setcounter{enumi}{1}
\tightlist
\item
  Let \(E_{0}\) be an infinite dimensional Banach space. Let E be the vector space that is the direct sum of \(E_{0}\) of \(\mathbf{R}^{(\mathbf{N})}\) and let \(E_{p}\) be the subspace of E that is the direct sum of \(E_{0}\) and of \(R^{p}\) (identified as the product of the first p factors of \(\mathbf{R}^{(\mathbf{N})}\) ; we give to \(E_{p}\) the product topology of those of its factor, so that the topology of \(E_{p}\) is induced by that of \(E_{p+1}\) . Show that on E the inductive limit topology of those of the \(E_{p}\) is not the finest of the topologies (compatible or not with the vector-space structure of E) which induces on each \(E_{p}\) a coarser topology than that of \(E_{p}\) . We can proceed as follows :
\end{enumerate}

\(\alpha)\) Let \(q\) be a norm on \(\mathbf{E}_0\) which defines a topology strictly coarser than that of \(\mathbf{E}_0\) (II, p.~74, exerc. 5). For every \(\varepsilon > 0\) define a mapping \(f_{\varepsilon}\) of \(\mathbf{E}_0\) in \(\mathbf{R}_+\) by the relation \(f_{\varepsilon}(x) = \sup(q(x), \varepsilon - \| x \|)\) . Show that \(f_{\varepsilon}\) is continuous and \(>0\) in \(\mathbf{E}_0\) and that

\[
\inf_{\| x \| = \varepsilon} f_{\varepsilon}(x) = 0 .
\]

\(\beta)\) Let \(U\) be the subset of \(E\) formed by the \((x, (t_n))\) such that \(t_n < f_{1/n}(x)\) for all \(n\) . Show that \(U \cap E_p\) is open in \(E_p\) for all \(p\) .

\(\gamma)\) Show that if \(\mathrm{V} \subset \mathrm{U}\) is an absorbent convex set, then \(\mathrm{V} \cap \mathrm{E}_0\) cannot contain any ball with centre 0 in \(\mathrm{E}_0\) .

\begin{enumerate}
\def\labelenumi{\arabic{enumi})}
\setcounter{enumi}{16}
\tightlist
\item
  For a subset A of a commutative group G, written additively, and for each \(n > 0\) denote the set of elements of the form \(\sum_{i=1}^{n} x_i\) , where \(x_i \in \mathbf{A}\) for all \(i\) by \({}^{n}+\mathbf{A}\) . We say that the set A
\end{enumerate}

of G is convex if, for every integer n > 0 the relation \(nx \in {}^{n}+A\) implies \(x \in A\) .

\begin{enumerate}
\def\labelenumi{\alph{enumi})}
\item
  Show that if a commutative topological group G (written additively) is isomorphic to a subgroup of the additive group of a locally convex vector space (with the induced topology) then there exists a fundamental system of symmetric convex neighbourhoods of 0 in G.
\item
  Conversely, let G be a Hausdorff topological commutative group (written additively) in which there exists a fundamental system, B, of symmetric convex neighbourhoods of 0. Show that G is without torsion, and, hence, can be considered (algebraically) as a subgroup of the additive group of a vector space on the field Q (A, II, § 7.10, cor. 1 to prop. 26). For every set V ∈ B, let \(\tilde{V}\) be the set of elements rx where x ∈ V and r varies in the set of rational numbers such that \(0 \leqslant r \leqslant 1\) ; show that \(\tilde{V}\) is symmetric and convex (in the sense defined above). Deduce, further, that if there is no open subgroup of G distinct from G itself, then the sets \(\tilde{V}\) form a fundamental system of neighbourhoods of 0 for a topology compatible with the vector space structure of E on Q (Q being given its usual topology); conclude that in this case G is isomorphic to a subgroup of the additive group of a Hausdorff locally convex space.
\item
  Let G be the group \(R \times R\) ordered lexicographically (A, VI, p.~7); consider the Hausdorff topology \(\mathcal{T}_{0}(G)\) on G that is compatible with its group structure (GT, IV, § 1, exerc. 1). Show that for this topology there exists a fundamental system of symmetric convex neighbourhoods of 0, but that G is not isomorphic to any subgroup of the additive group of a Hausdorff topological vector space over R.
\end{enumerate}

